\begin{frame}
\frametitle{Instrukce lw -- load word}

\textbf{\texttt{lw -- load word}} načti slovo z paměti do registru

\begin{tabular}{|l|l|}\hline
Popis & Načte slovo do registru ze zadané adresy z paměti \\ \hline
Operace & \texttt{[rd]} $\leftarrow$ \texttt{Mem[[rs1]+imm12]} \\ \hline
Syntax & lw rd, imm12(rs1) \\ \hline
Kódování & \texttt{iiii iiii iiii ssss s010 dddd d000 0011} \\ \hline
 & s -- rs1; d -- rd; i -- konstanta \\ \hline
\end{tabular}

\bigskip

Příklad: Načtěte 32-bitové slovo z paměti z adresy 0x400 do registru 2:\\
\textbf{\texttt{lw x2, 0x400(x0)}}

\textbf{\texttt{\color{red}iiii iiii iiii}}\phantom{xx}\textbf{\texttt{\color{blue}ssss s}}\hspace{0.1cm}\textbf{\texttt{010\hspace{0.25cm}{\color{green}dddd d}\hspace{0.05cm}000 0011}}\\
$\underbrace{\textbf{\texttt{\color{red}0100 0000 0000}}}_{0x400}\texttt{ }\underbrace{\textbf{\texttt{\color{blue}0000 0}}}_{0}\underbrace{\textbf{\texttt{010}}}_{func3}\phantom{i}\underbrace{\textbf{\texttt{\color{green}0001 0}}}_{2}\underbrace{\textbf{\texttt{000 0011}}}_{opcode}$\\

\textbf{\texttt{0x 40 00 21 03}} -- kód pro instrukci \textbf{\texttt{lw x2, 0x400(x0)}}

Poznámka: Registr x0 má natvrdo hodnotu 0, kterou nelze změnit

\end{frame}

\begin{frame}[shrink=15]
\frametitle{Implementace instrukce lw}

\textbf{\texttt{lw}}: \texttt{rs1} -- základní adresa, \texttt{imm12} -- posunutí adresy, \texttt{rd} -- registr, kam se uloží data z paměti

\bigskip

\begin{table}
\footnotesize
\begin{tabular}{|m{0.4cm}|m{0.4cm}|m{1.0cm}|m{1.0cm}|m{0.4cm}|m{1.0cm}|m{1.0cm}|m{1.0cm}|m{0.4cm}|m{1.0cm}|}\hline
Typ & 31 & 30...25 & 24...21 & 20 & 19...15 & 14...12 & 11...8 & 7 & 6...0 \\ \hline
I & \multicolumn{4}{c|}{ imm[11:0] } & rs1 & fnct3 &\multicolumn{2}{c|}{ rd } & opcode\\ \hline
\end{tabular}
\end{table}

\bigskip

\begin{center}
\includegraphics[width=0.85\textwidth]{generated/cpu/instruction/cpu-load.pdf}
\end{center}

\end{frame}

\begin{frame}[shrink=15]
\frametitle{Implementace instrukce lw}

\textbf{\texttt{lw}}: \texttt{rs1} -- základní adresa, \texttt{imm12} -- posunutí adresy, \texttt{rd} -- registr, kam se uloží data z paměti

\bigskip

\begin{table}
\footnotesize
\begin{tabular}{|m{0.4cm}|m{0.4cm}|m{1.0cm}|m{1.0cm}|m{0.4cm}|m{1.0cm}|m{1.0cm}|m{1.0cm}|m{0.4cm}|m{1.0cm}|}\hline
Typ & 31 & 30...25 & 24...21 & 20 & 19...15 & 14...12 & 11...8 & 7 & 6...0 \\ \hline
I & \multicolumn{4}{c|}{ imm[11:0] } & rs1 & fnct3 &\multicolumn{2}{c|}{ rd } & opcode\\ \hline
\end{tabular}
\end{table}

Zápis proběhna na rostoucí hranu hodin.

\begin{center}
\includegraphics[width=0.90\textwidth]{generated/cpu/instruction/cpu-load-red.pdf}
\end{center}



\end{frame}


\begin{frame}[shrink=15]
\frametitle{Implementace instrukce lw}

\textbf{\texttt{lw}}: \texttt{rs1} -- základní adresa, \texttt{imm12} -- posunutí adresy, \texttt{rd} -- registr, kam se uloží data z paměti

\bigskip

\begin{table}
\footnotesize
\begin{tabular}{|m{0.4cm}|m{0.4cm}|m{1.0cm}|m{1.0cm}|m{0.4cm}|m{1.0cm}|m{1.0cm}|m{1.0cm}|m{0.4cm}|m{1.0cm}|}\hline
Typ & 31 & 30...25 & 24...21 & 20 & 19...15 & 14...12 & 11...8 & 7 & 6...0 \\ \hline
I & \multicolumn{4}{c|}{ imm[11:0] } & rs1 & fnct3 &\multicolumn{2}{c|}{ rd } & opcode\\ \hline
\end{tabular}
\end{table}

\bigskip

\begin{center}
\includegraphics[width=0.88\textwidth]{generated/cpu/instruction/cpu-load2.pdf}
\end{center}
\end{frame}


\begin{frame}
\frametitle{QtRvSim - RISC-V Simulátor}

\includegraphics[width=0.9\textwidth]{generated/simulator/simulator.pdf}

\end{frame}



\begin{frame}
\frametitle{Instrukce sw -- store word}

\textbf{\texttt{sw -- store word}} ulož slovo z registru do paměti

\bigskip

\begin{tabular}{|l|l|}\hline
Popis & Uloží slovo z registru rs2 do zadané adresy z paměti \\ \hline
Operace & \texttt{Mem[[rs1]+imm12]} $\leftarrow$ \texttt{[rs2]} \\ \hline
Syntax & sw rs2, imm12(rs1) \\ \hline
Kódování & \texttt{iiii iiit tttt ssss s010 iiii i010 0011} \\ 
 & t -- rs2; s -- rs1; i -- konstanta \\ \hline
\end{tabular}

\bigskip

Příklad: Uloží 32-bitové slovo z registru 2 do paměti na adresu 0x404 plus hodnota registru 5:\\
\textbf{\texttt{sw x2, 0x404(x5)}}

\textbf{\texttt{{\color{green}iiii iii}\hspace{0.08cm}\color{red}t tttt}}\phantom{x}\hspace{0.13cm}\textbf{\texttt{\color{blue}ssss s}}\hspace{0.1cm}\textbf{\texttt{010\hspace{0.25cm}{\color{green}iiii i}\hspace{0.05cm}010 0011}}\\
$\underbrace{\textbf{\texttt{\color{green}0100 000}}}_{0x40\_}
\underbrace{\textbf{\texttt{\color{red}0 0010}}}_{2}
\texttt{ }\underbrace{\textbf{\texttt{\color{blue}0010 1}}}_{5}
\underbrace{\textbf{\texttt{010}}}_{func3}\phantom{i}\underbrace{\textbf{\texttt{\color{green}0010 0}}}_{0x \_04}
\underbrace{\textbf{\texttt{010 0011}}}_{opcode}$\\

\textbf{\texttt{0x 40 22 a2 23}} -- kód pro instrukci \textbf{\texttt{sw x2, 0x404(x5)}}

\end{frame}

\begin{frame}[shrink=18]
\frametitle{Implementace instrukce sw}

\textbf{\texttt{sw}}: \texttt{rs1} -- základní adresa, \texttt{imm12} -- posunutí adresy, \texttt{rs2} -- registr, který se uloží do paměti

\bigskip

\begin{table}
\footnotesize
\begin{tabular}{|m{0.4cm}|m{0.4cm}|m{1.0cm}|m{1.0cm}|m{0.4cm}|m{1.0cm}|m{1.0cm}|m{1.0cm}|m{0.4cm}|m{1.0cm}|}\hline
Typ & 31 & 30...25 & 24...21 & 20 & 19...15 & 14...12 & 11...8 & 7 & 6...0 \\ \hline
S & \multicolumn{2}{c|}{ imm[11:5] } & \multicolumn{2}{c|}{ rs2 } & rs1 & fnct3 &\multicolumn{2}{c|}{ imm[4:0] } & opcode\\ \hline
\end{tabular}
\end{table}

\bigskip

\includegraphics[width=0.90\textwidth]{generated/cpu/instruction/cpu-store2.pdf}

\end{frame}




\begin{frame}
\frametitle{Instrukce add -- sečti dva registry}

\textbf{\texttt{add -- addition}} -- sečti hodnoty dvou registrů a součet ulož do registru

\bigskip

\begin{tabular}{|l|l|}\hline
Popis & Sečti hodnoty z registrů rs1 a rs2 a výsledek ulož do rd \\ \hline
Operace & \texttt{[rd]} $\leftarrow$ \texttt{[rs1] + [rs2]} \\ \hline
Syntax & add rd, rs1, rs2 \\ \hline
Kódování & \texttt{0000 000t tttt ssss s000 dddd d011 0011} \\
 & t -- rs2; s -- rs1; d -- rd \\ \hline
\end{tabular}

\bigskip

Příklad: Sečti hodnoty z registru 2 a registru 3, výsledek ulož do registru 4:\\
\textbf{\texttt{add x4, x2, x3}}

\textbf{\texttt{0000 000\hspace{0.08cm}\color{red}t tttt}}\phantom{x}\hspace{0.13cm}\textbf{\texttt{\color{blue}ssss s}}\hspace{0.1cm}\textbf{\texttt{000\hspace{0.25cm}{\color{green}dddd d}\hspace{0.05cm}011 0011}}\\
$\underbrace{\textbf{\texttt{0000 000}}}_{funct7}
\underbrace{\textbf{\texttt{\color{red}0 0011}}}_{3}
\texttt{ }\underbrace{\textbf{\texttt{\color{blue}0001 0}}}_{2}
\underbrace{\textbf{\texttt{000}}}_{func3}\phantom{i}
\underbrace{\textbf{\texttt{\color{green}0010 0}}}_{4}
\underbrace{\textbf{\texttt{011 0011}}}_{opcode}$\\

\textbf{\texttt{0x 00 31 02 33}} -- kód pro instrukci \textbf{\texttt{add x4, x2, x3}}


\end{frame}

\begin{frame}[shrink=25]
\frametitle{Implementace instrukce add}

\textbf{\texttt{add}}: \texttt{rs1}, \texttt{rs2} -- sčítance, \texttt{rd} -- registr, kam se uloží součet

\bigskip

\begin{table}
\footnotesize
\begin{tabular}{|m{0.4cm}|m{0.4cm}|m{1.0cm}|m{1.0cm}|m{0.4cm}|m{1.0cm}|m{1.0cm}|m{1.0cm}|m{0.4cm}|m{1.0cm}|}\hline
Typ & 31 & 30...25 & 24...21 & 20 & 19...15 & 14...12 & 11...8 & 7 & 6...0 \\ \hline
R & \multicolumn{2}{c|}{ fnct7 } & \multicolumn{2}{c|}{ rs2 } & rs1 & fnct3 &\multicolumn{2}{c|}{ rd } & opcode\\ \hline
\end{tabular}
\end{table}

\bigskip

\includegraphics[width=0.85\textwidth]{generated/cpu/instruction/cpu-add.pdf}

\end{frame}


\begin{frame}
\frametitle{Implementace instrukce sub, and, or, slt}

Jediný rozdíl oproti instrukci sčítání add je v hodnotě ALUControl, který vybírá, jakou operaci jednotka ALU provede. Cesta dat a uložení výsledku je stejná jako u sčítání.
\bigskip

\bigskip

\includegraphics[width=0.85\textwidth]{generated/cpu/instruction/cpu-add2.pdf}

\end{frame}


\begin{frame}
\frametitle{Instrukce addi -- přičti konstantu k registru}

\textbf{\texttt{addi -- addition immediate}} -- sečti hodnotu registru a konstanty a výsledek ulož do registru

\bigskip

\begin{tabular}{|l|l|}\hline
Popis & Sečti hodnoty z registru rs1 a imm12 a výsledek ulož do rd \\ \hline
Operace & \texttt{[rd]} $\leftarrow$ \texttt{[rs1] + imm12} \\ \hline
Syntax & addi rd, rs1, imm12 \\ \hline
Kódování & \texttt{iiii iiii iiii ssss s000 dddd d001 0011} \\ 
 & i -- konstanta; s -- rs1; d -- rd \\ \hline
\end{tabular}

\bigskip

Příklad: Zvětši hodnotu registru 7 o 4 (výsledek tedy ulož do registru 7):\\
\textbf{\texttt{addi x7, x7, 4}}

\textbf{\texttt{\color{red}iiii iiii iiii}}\phantom{xx}\textbf{\texttt{\color{blue}ssss s}}\hspace{0.1cm}\textbf{\texttt{000\hspace{0.25cm}{\color{green}dddd d}\hspace{0.05cm}001 0011}}\\
$\underbrace{\textbf{\texttt{\color{red}0000 0000 0100}}}_{0x004}\texttt{ }\underbrace{\textbf{\texttt{\color{blue}0011 1}}}_{7}\underbrace{\textbf{\texttt{000}}}_{func3}\phantom{i}\underbrace{\textbf{\texttt{\color{green}0011 1}}}_{7}\underbrace{\textbf{\texttt{001 0011}}}_{opcode}$\\

\textbf{\texttt{0x 00 43 83 93}} -- kód pro instrukci \textbf{\texttt{addi x7, x7, 4}}


\end{frame}


\begin{frame}[shrink=10]
\frametitle{Implementace instrukcí addi, ori, andi}

\textbf{\texttt{addi -- add immediate}}: \texttt{[rd]} $\leftarrow$ \texttt{[rs1] + imm12} -- přičte k hodnotě registru \texttt{rs1} konkrétní číslo \texttt{imm} a výsledek uloží do registru \texttt{rd}

\bigskip

\begin{table}
\footnotesize
\begin{tabular}{|m{0.4cm}|m{0.4cm}|m{1.0cm}|m{1.0cm}|m{0.4cm}|m{1.0cm}|m{1.0cm}|m{1.0cm}|m{0.4cm}|m{1.0cm}|}\hline
Typ & 31 & 30...25 & 24...21 & 20 & 19...15 & 14...12 & 11...8 & 7 & 6...0 \\ \hline
I & \multicolumn{4}{c|}{ imm[11:0] } & rs1 & fnct3 &\multicolumn{2}{c|}{ rd } & opcode\\ \hline
\end{tabular}
\end{table}

\bigskip

\includegraphics[width=0.85\textwidth]{generated/cpu/instruction/cpu-addi.pdf}

\end{frame}


\begin{frame}[shrink=15]
\frametitle{Implementace instrukce skoku beq}

\textbf{\texttt{beq -- branch if equal}}: \texttt{[pc]} $\leftarrow$ \texttt{[pc] + SignImm} -- přičte k hodnotě registru \texttt{pc} konkrétní znaménkové číslo rozšířené na 32 bitů \texttt{imm} a výsledek uloží do registru \texttt{pc} pokud nastala rovnost hodnot v registrech rs1 a rs2.

\bigskip

\begin{table}
\footnotesize
\begin{tabular}{|m{0.4cm}|m{0.4cm}|m{1.0cm}|m{1.0cm}|m{0.4cm}|m{1.0cm}|m{1.0cm}|m{1.0cm}|m{0.4cm}|m{1.0cm}|}\hline
Typ & 31 & 30...25 & 24...21 & 20 & 19...15 & 14...12 & 11...8 & 7 & 6...0 \\ \hline
B & imm [12] & imm [10:5]  & \multicolumn{2}{c|}{ rs2 } & rs1 & fnct3 & imm[4:1]& imm [11] & opcode\\ \hline
\end{tabular}
\end{table}

\bigskip

\includegraphics[width=0.85\textwidth]{generated/cpu/instruction/cpu-beq.pdf}

\end{frame}
